\section{Monitoreo, observabilidad y gobernanza}

\subsection{Señales de monitoreo de despliegue}

\begin{table}[H]
\centering
\caption{Lista de verificación de monitoreo de Deployment}
\begin{tabular}{p{0.25\textwidth}p{0.28\textwidth}p{0.25\textwidth}}
\toprule
Señal & Significado & Acción activada \\
\midrule
Latency increase & Degradación del desempeño del modelo o la herramienta & Revertir al last-known-good prompt y agregar una optimizer sample \\
Tool error & Respuesta anómala de la herramienta & Suspender el módulo correspondiente, emitir una alerta y hacer una verificación secundaria \\
Metric drift & Caída en la puntuación & Volver a recopilar los datos de entrenamiento e iniciar el auto-optimizer \\
User feedback spike & Aumento súbito en la tasa de rechazo de respuesta o de malentendidos & Incorporar una revisión humana y actualizar el evidence trace \\
\bottomrule
\end{tabular}
\end{table}

Estas señales normalmente son ingested al observability stack. Latency y Tool error hacen que el alerting se active de inmediato; Metric drift, en cambio, necesita schedule una ejecución comparativa (por ejemplo, un nightly run). User feedback spike normalmente usa QA note + Slack channel (como `\#agent-monitors`) para hacer triage manual.

\begin{warningbox}{Ciclo cerrado de gobernanza}
En cada despliegue importante hay que ejecutar el three-step workflow: define → monitor → annotate. El paso annotate es fundamental, pues asegura que el optimizer posterior siga recordando por qué se configuró así en un inicio.
\end{warningbox}

\subsection{Action Timeline}

En el README de cada release hay que actualizar la tabla de Action Timeline, donde se enumeran el nombre del módulo, la condición que lo activa, el metric threshold y el responsible owner. De esta manera, cuando surge un problema es posible ubicar rápidamente la cadena de responsabilidad.

Action Timeline también puede vincularse con la firma del colaborador: cada ajuste del pipeline debe llevar ``who changed what'', lo que permite deslindar responsabilidades rápidamente tras un incidente. Para equipos multiinquilino, se recomienda guardar también cada timeline entry en el central Incident Log, y enlazarlo mediante el `Evidence Matrix ID`.

\subsection{Resumen de la sección}

El monitoreo y la gobernanza son la tercera capa de protección del DSPy pipeline: es necesario convertir señales como latency, metric y feedback en respuestas ejecutables, en lugar de dejarlas sin control.

